\section{Numerical Results}

In this section, we present numerical experiments illustrating the performance of the proposed 
finite element approximation and optimal control framework.

All computations are intended to be carried out using the FEniCSx finite element library. 
The state equation is solved forward in time, while the adjoint equation is solved backward 
in time. The control variable is updated through a gradient-based optimization algorithm.

\subsection{Test Problem}

We consider the square domain

\[
D=(0,1)^2,
\]

and the final time

\[
T=1.
\]

The diffusion coefficient is chosen as

\[
\kappa = 10^{-2},
\]

while the velocity field is defined by

\[
\mathbf b(x,y)
=
\begin{pmatrix}
1 \
0
\end{pmatrix}.
\]

The desired state is given by

\[
y_d(x,y)
=
\sin(\pi x)\sin(\pi y).
\]

The initial condition is taken as

\[
y_0(x,y)=0.
\]

Homogeneous Dirichlet boundary conditions are imposed on the whole boundary of the domain.

The regularization parameter is chosen as

\[
\alpha =10^{-3}.
\]

\subsection{Finite Element Parameters}

The computational domain is discretized using a structured triangular mesh.

Continuous piecewise linear finite elements \((P_1)\) are employed for the approximation of the state, 
adjoint, and control variables.

The time interval is divided into \(N_t\) uniform time steps with time increment

\[
\Delta t
=
\frac{T}{N_t}.
\]

The state equation is discretized using the backward Euler method, while the optimization procedure 
relies on a gradient descent algorithm.

\subsection{State Solution}

The first experiment consists in solving the state equation for a prescribed control.

Figure~\ref{fig:state_solution} illustrates the numerical approximation of the state variable at the 
final time \(T\).

\begin{figure}[H]
\centering
\includegraphics[width=0.65\textwidth]{figures/state/state_solution.png}
\caption{Numerical approximation of the state variable at the final time.}
\label{fig:state_solution}
\end{figure}

The solution exhibits the expected advection and diffusion effects induced by the transport field 
and the diffusion coefficient.

\subsection{Adjoint Solution}

The adjoint equation is solved backward in time starting from the terminal condition

\[
p(T)=0.
\]

Figure~\ref{fig:adjoint_solution} presents the numerical approximation of the adjoint variable.

\begin{figure}[H]
\centering
\includegraphics[width=0.65\textwidth]{figures/state/adjoint_solution.png}
\caption{Numerical approximation of the adjoint variable.}
\label{fig:adjoint_solution}
\end{figure}

The adjoint solution provides sensitivity information required for the optimization process.

\subsection{Optimal Control}

Using the optimality condition

\[
u^\ast
=
\frac{1}{\alpha}p^\ast,
\]

the optimal control is computed iteratively.

Figure~\ref{fig:optimal_control} displays the resulting control distribution.

\begin{figure}[H]
\centering
\includegraphics[width=0.65\textwidth]{figures/control/optimal_control.png}
\caption{Optimal control obtained after convergence of the optimization algorithm.}
\label{fig:optimal_control}
\end{figure}

The computed control successfully drives the state toward the prescribed target state while 
maintaining a moderate control cost.

\subsection{Convergence of the Optimization Algorithm}

To assess the convergence of the optimization procedure, we monitor the evolution of the cost 
functional during the iterations.

Figure~\ref{fig:cost_history} shows a typical convergence history.

\begin{figure}[H]
\centering
\includegraphics[width=0.65\textwidth]{figures/control/cost_history.png}
\caption{Evolution of the cost functional during the optimization process.}
\label{fig:cost_history}
\end{figure}

A monotonic decrease of the objective functional is observed, indicating the effectiveness of 
the gradient-based optimization algorithm.

\subsection{Influence of the Noise Intensity}

We next investigate the influence of the stochastic forcing parameter \(\sigma\).

Several simulations are performed for different values of \(\sigma\).

The results indicate that increasing the noise intensity leads to larger fluctuations in the state 
variable and generally requires stronger control actions to achieve the desired target state.

\subsection{Discussion}

The numerical experiments confirm the theoretical results established in the previous sections.

The finite element approximation provides stable and accurate numerical solutions for both the 
state and adjoint equations. Furthermore, the gradient-based optimization strategy successfully 
computes controls that minimize the objective functional.

These results demonstrate the applicability of the proposed framework to stochastic transport 
phenomena arising in environmental and climate applications.
